% Frozen prospective decision criteria for the Materna analysis.
\section*{Prospectively Frozen Decision Criteria}

All resource-classification conditions are conjunctive.  A positive horizon,
VM, fold, or seed aggregate has gain strictly greater than zero.  A positive
resource is classified as either \textsc{Broad Positive} or
\textsc{Conditionally Positive}.

\begin{description}
  \item[\textsc{Broad Positive}.] Overall gain is at least $+1\%$; at least
    $4/5$ horizons, $25/35$ VMs, $4/5$ folds, and $2/3$ seed aggregates are
    positive.
  \item[\textsc{Conditionally Positive}.] The resource is not \textsc{Broad
    Positive}; overall gain is strictly positive; and at least $3/5$ horizons,
    $15/35$ VMs, $3/5$ folds, and $2/3$ seed aggregates are positive.
  \item[\textsc{Broad Negative}.] Overall gain is at most $-1\%$; and at most
    $1/5$ horizons, $10/35$ VMs, $1/5$ folds, and $1/3$ seed aggregates are
    positive.
  \item[\textsc{Neutral or Mixed}.] Every case not assigned to one of the three
    classes above.
\end{description}

The equal-temporal-block synthesis is the unweighted arithmetic mean of the T1,
T2, and T3 primary gains.  A materially negative resource has overall gain at
most $-1\%$.  Strong target-specific trade-offs exist when the maximum
resource-level gain is at least $+1\%$ and the minimum is at most $-1\%$.
Material positive-and-negative synthesis exists when at least one synthesised
resource gain is at least $+1\%$ and another is at most $-1\%$.

\subsection*{Trace-level comparison decisions}

For S2 versus S1, \textsc{Cross-Resource Input Supported} requires at least
three positive resources, positive equal-resource mean primary gain, and at
most one \textsc{Broad Negative} resource.  \textsc{Cross-Resource Input Not
Supported} applies when at most one resource is positive and the equal-resource
mean is non-positive, or when at least two resources are \textsc{Broad
Negative}.  All other cases are \textsc{Cross-Resource Input Mixed}.

For M1 versus S2, conservative precedence is: \textsc{Negative Transfer},
\textsc{Not Supported}, \textsc{Supported}, then \textsc{Conditionally
Supported}.  \textsc{Supported} requires at least three positive resources,
positive equal-resource mean primary gain, and no \textsc{Broad Negative}
resource.  \textsc{Conditionally Supported} requires a positive equal-resource
mean and at least two positive resources when the full support criterion fails.
\textsc{Negative Transfer} applies when at least one resource is positive and
at least two resources are either \textsc{Broad Negative} or materially
negative, or when strong target-specific trade-offs occur.  \textsc{Not
Supported} applies when at most one resource is positive and the equal-resource
mean is non-positive; a residual case with no named predicate is also assigned
\textsc{Not Supported} with the recorded residual-no-support reason.

For M1 versus S1, \textsc{Joint System Broadly Beneficial} requires positive
overall gain for at least three resources and a positive equal-resource mean.
\textsc{Joint System Broadly Harmful} requires negative overall gain for at
least three resources and a negative equal-resource mean.  All other cases are
\textsc{Joint System Mixed}.

\subsection*{RQ1--RQ4 synthesis decisions}

RQ1 uses the following precedence.  \textsc{Broad Cross-Resource Benefit}
requires supported decisions in at least two traces, a positive mean of the
four synthesised resource gains, and at least three positive synthesised
resources.  If broad benefit fails, \textsc{Resource--Temporal Dependence
Confirmed} applies when trace-level cross-resource decisions are not identical,
at least two resources change sign, or at least two resources change
classification.  \textsc{Cross-Resource Input Not Supported} requires zero
supported traces, a non-positive four-resource synthesis mean, at most one
positive synthesised resource, and no temporal-dependence decision.  Otherwise
the decision is \textsc{Mixed or Inconclusive}.

RQ2 uses conservative precedence.  \textsc{Negative Sharing Replicated}
requires zero supported traces, at most one conditionally supported trace, a
non-positive four-resource synthesis mean, and at most one positive synthesised
resource.  \textsc{Mixed Sharing Effect} requires zero supported traces and
material positive-and-negative synthesis.  \textsc{Sharing Benefit Replicated}
requires at least two supported or conditionally supported traces, a positive
four-resource synthesis mean, and at least three positive synthesised
resources.  Otherwise the decision is \textsc{Inconclusive}.

RQ3 uses the following precedence.  \textsc{Joint System Underperformance
Replicated} requires all four synthesised resource gains to be negative and at
least two broadly harmful trace decisions.  \textsc{Joint System Partial
Underperformance} applies when the four-resource synthesis mean is non-positive,
at most two synthesised resources are positive, and the full replicated
criterion fails.  \textsc{Joint System Benefit} requires at least three
positive synthesised resources, a positive four-resource synthesis mean, and at
least two broadly beneficial trace decisions.  Otherwise the decision is
\textsc{Mixed or Inconclusive}.

For RQ4, a resource--comparison pair is temporally classification-variable when
its frozen classification is not identical across T1, T2, and T3; raw sign
reversals are reported separately.  \textsc{Temporal Heterogeneity Confirmed}
applies when at least two resource--comparison pairs are classification-variable
or when the global trace decisions differ for S2 versus S1 or M1 versus S2.
\textsc{Temporal Heterogeneity Weak} applies when exactly one pair is
classification-variable and global trace decisions are identical.
\textsc{Temporal Heterogeneity Not Supported} applies when no pair is
classification-variable and global trace decisions are identical.  Post-hoc
regime mining is not permitted.
